% ============================================================
% 職務経歴書（通常組版版）— jlreq 両端揃え
% 元: 職務経歴書_陳凱_20260829.txt / build.tex（Verbatim版）
% 使い方: docker run --rm -v "$PWD":/workdir paperist/texlive-ja:latest latexmk -lualatex build_jlreq.tex
% ============================================================
\documentclass[paper=a4paper,fontsize=10pt,head_space=18mm,foot_space=15mm,gutter=15mm,fore_edge=15mm,baselineskip=16pt]{jlreq}

\usepackage{luatexja}   % 日本語
\usepackage[top=18mm,bottom=18mm,left=16mm,right=16mm]{geometry}
\usepackage{array}
\usepackage{booktabs}
\usepackage{xcolor}
\usepackage[hidelinks]{hyperref}

\setlength{\parindent}{0pt}
\setlength{\itemsep}{1pt}
\setlength{\parsep}{0pt}
\renewcommand{\baselinestretch}{1.15}

% セクション見出し（■相当）
\newcommand{\kyariku}[1]{%
  \par\vspace{10pt}%
  {\Large\bfseries #1}\par\vspace{4pt}%
  {\color{gray}\rule{\textwidth}{0.4pt}}\par\vspace{6pt}%
}

% プロジェクト見出し
\newcommand{\project}[1]{%
  \par\vspace{6pt}%
  {\large\bfseries #1}\par\vspace{2pt}%
}


\begin{document}

% ============ ヘッダ ============
\begin{center}
  {\Huge\bfseries 職\ 務\ 経\ 歴\ 書}\\[6pt]
  2026年9月4日\\[2pt]
  陳\ 凱
\end{center}

\vspace{4pt}

% ============ 職務要約 ============
\kyariku{職務要約}

2019年より本国の義務兵役に従事し、所属機関にてサーバー稼働監視システムの開発を約2年3ヶ月担当。C\#/WPF によるデスクトップクライアントと ASP.NET Core による API サーバーの双方を一人で設計・実装いたしました。その後、RPA製品を開発する企業にて業務自動化ツールのUI開発・改修に従事し、大学へ編入学。2024年に卒業。

2026年に来日し、現在は損害保険会社向け基幹システムをオンプレミスから AWS へ移行する案件（大手SIer元請・12名体制）に開発メンバーとして参画しております。既存システムのソースコード解析から外部設計書・内部設計書の作成、単体テストまでの上流〜下流工程を一貫して担当しております。

移行先環境を理解する必要があると考え、AWS認定資格を CLF・SAA・SAP まで取得。業務外でも個人開発を継続し、設計判断を ADR として文書化する進め方を実践しております。日本語は JLPT N1、中国語はネイティブです。

% ============ 活かせる経験・スキル ============
\kyariku{活かせる経験・スキル}

\begin{itemize}
  \item 金融（損害保険）業界の基幹システム開発経験
  \item オンプレミス環境から AWS への移行案件における設計工程の経験
  \item レガシーシステムのソースコード解析による現行仕様の復元
  \item 日本語での外部設計書・内部設計書の作成
  \item サーバー稼働監視システムの設計・開発（リソース監視・異常検知・アラート通知）
  \item エンドユーザー向け業務アプリケーションのUI設計・改修（C\# / WPF / MVVM）
  \item フロントエンド／バックエンド双方の実装経験（React・TypeScript ／ C\#・Java）
  \item AWSアーキテクチャに関する体系的知識（AWS認定 SAP 保有）
  \item 日本語による業務遂行能力（JLPT N1）／中国語ネイティブ
\end{itemize}

% ============ テクニカルスキル ============
\kyariku{テクニカルスキル}

\begin{tabular}{@{}p{4.5em}p{\dimexpr\textwidth-4.5em\relax}@{}}
  言語 & Java（現職・約6ヶ月）／ \textbf{C\#（実務約3年）} ／ TypeScript・React（2023年〜・個人開発含む）／ Python（個人開発）／ SQL \\
  フレームワーク & Thymeleaf ／ ASP.NET Core ／ Entity Framework Core ／ WPF・Prism (MVVM) ／ React 19 ／ FastAPI \\
  クラウド & AWS（EC2・認定資格 CLF / SAA / SAP） \\
  その他 & Docker ／ Git ／ SQLite \\
\end{tabular}

\vspace{2pt}
{\small ※言語の使用年数は実務を基準とし、個人開発での継続利用は併記。}

% ============ 職務経歴 ============
\kyariku{職務経歴}

\project{【1】株式会社日東ソフト（2026年4月 〜 現在）}
事業内容：システム開発・技術者派遣／雇用形態：契約社員（2026年11月30日 契約期間満了予定）

\project{■ 損害保険会社向け 申込管理システム クラウド移行案件}
期間：2026年4月 〜 現在（進行中）／概要：大手SIer元請の金融系基幹システム刷新案件／規模：12名（開発チーム所属・別途インフラ基盤チーム）／役割：開発メンバー

\textbf{▼ プロジェクト概要}\par
長年オンプレミス環境で運用されてきた損害保険の申込管理システムを、AWS環境へ移行するリプレイス案件。レガシーな Java フレームワークで構築された既存システムを Thymeleaf ベースの構成へ刷新し、アプリケーションサーバを AWS EC2 上へ移行する。

\textbf{▼ 担当業務}
\begin{itemize}
  \item 既存システムのソースコード解析および現行仕様の調査
  \item 外部設計書・内部設計書の作成（日本語）
  \item 移行方式に関する設計レビューへの参加
  \item 単体テスト（UT）仕様書の作成およびテスト実施
  \item 移行前後の動作差分の検証
\end{itemize}

\textbf{▼ 使用技術}\par
Java / Thymeleaf / AWS（EC2）/ Git

\textbf{▼ 本案件を通じて}\par
ドキュメントが十分に整備されていないレガシーシステムに対し、ソースコードから現行仕様を復元した上で設計書を作成し、レビューを通過させました。またインフラ基盤チームが AWS 側を担当する体制の中で、開発側としても移行先環境の全体像を理解する必要があると考え、AWS認定資格（CLF / SAA / SAP）を取得。クラウド前提の設計における考慮点について、基盤チームと同じ論点で議論できる土台を作りました。

\project{【2】武漢慧霊控科技有限公司（2022年2月 〜 2022年8月）}
事業内容：RPA（Robotic Process Automation）製品の開発・提供

\project{■ RPA デスクトップアプリケーションのUI開発・改修}
期間：2022年2月 〜 2022年8月（7ヶ月）／役割：開発メンバー（UI領域を担当）

\textbf{▼ プロジェクト概要}\par
企業の定型業務を自動化する RPA 製品のデスクトップアプリケーション。利用者が業務フローを画面上で組み立て、実行状況を監視する操作環境を C\# / WPF により開発・改修する。エンジニア以外の一般利用者が日常的に操作するため、操作性の設計が品質に直結する案件であった。

\textbf{▼ 担当業務}
\begin{itemize}
  \item WPF / XAML による画面の実装および既存UIの改修
  \item MVVM パターンに基づく ViewModel の実装とデータバインディングの整備
  \item 利用者からの要望・問い合わせを起点とした操作性の改善
  \item 不具合の原因調査および修正（既存コードの解析を含む）
  \item 改修時の影響範囲の特定と、既存動作の非破壊性の確認
\end{itemize}

\textbf{▼ 使用技術}\par
C\# / WPF / XAML / .NET / MVVM / Git

\textbf{▼ 本案件を通じて}\par
既存の画面ロジックがコードビハインドに集中していた箇所について、改修の都度 ViewModel 側へ責務を移し、テスト・再利用が可能な状態へ段階的に整理しました。また「操作が分かりにくい」という抽象的な要望に対し、実際の操作手順を確認してどの工程で迷いが生じているかを特定し、画面遷移と表示情報の見直しとして具体化しました。大規模な既存コードベースに影響範囲を確認してから着手する進め方は、現在のレガシーシステム移行案件でも活きています。

\project{【3】公的機関（2019年9月 〜 2021年12月）※本国の義務兵役として従事}
※機関名および監視対象システムの詳細は、守秘義務の関係で記載を控えさせていただきます

\project{■ サーバー稼働監視・アラート通知システムの開発}
期間：2019年9月 〜 2021年12月（2年3ヶ月）／役割：設計・開発（クライアント／サーバー双方を担当）

\textbf{▼ プロジェクト概要}\par
複数のサーバーホストを対象とした稼働監視・アラート通知システム。デスクトップクライアントと、監視データを収集・蓄積する API サーバーを一人で設計・実装。各ホストの CPU 使用率・メモリ・ディスク空き容量に加え、システム間連携データの異常やファイル生成間隔を監視し、規定値を超過した場合にアラートを通知する。

\textbf{▼ 担当業務}
\begin{itemize}
  \item C\# / WPF によるデスクトップクライアントの設計・実装（Prism を用いた MVVM アーキテクチャの採用）
  \item ASP.NET Core / Entity Framework Core による REST API の設計・実装
  \item 監視データの収集・蓄積を行うデータベース設計
  \item 利用部門へのヒアリングによる監視要件の整理
\end{itemize}

\textbf{▼ 使用技術}\par
C\# / WPF / Prism (MVVM) / ASP.NET Core / Entity Framework Core / SQL

\textbf{▼ 本案件を通じて}\par
個別に確認していた各ホストの状態を一画面で把握できる運用に変更し、業務要件に基づく独自の異常検知ルールにより障害の早期検知に寄与しました。監視・運用の観点は、現在のクラウド移行案件においても設計上の判断材料となっています。

% ============ その他の開発経験 ============
\kyariku{その他の開発経験}

\project{【卒業研究】Web小説管理システム（2024年2月 〜 2024年4月）}
フロントエンドとバックエンドを分離した構成の Web アプリケーションを個人で設計・実装。
\begin{itemize}
  \item バックエンド：C\# / .NET による REST API の設計・実装
  \item フロントエンド：TypeScript / React による SPA の実装
  \item Docker によるコンテナ化と実行環境の構築
\end{itemize}
サーバーサイドとクライアントサイドの双方を一人で実装したことで、API 設計における責務の分担や、フロント・バックエンド間のインターフェース設計について実践的に理解する機会となりました。

※2023年以降、React / TypeScript を用いた Web アプリケーションの開発を継続しております。

% ============ 個人開発 ============
\kyariku{個人開発}

業務外で継続的に個人開発を行っております。実装には生成AIを併用しておりますが、アーキテクチャの選定と設計判断は自身で行い、その理由と却下した代替案を ADR（Architecture Decision Record）として文書化しています。

\textbf{Kotoba Studio — 日本語学習支援アプリケーション}\par
字幕のない日本語動画に単語レベルのタイムスタンプ付き字幕を生成し、動画・字幕・タイムラインを同期させた精読環境を提供するローカルアプリ。Python 3.11 / FastAPI / WhisperX / SQLite / React 19 / TypeScript\\
\url{https://github.com/RainYaya/bikuri-NihongoLens}

\textbf{baku-yomi（爆読）— 日中バイリンガル読書アプリ}\par
日中対訳EPUBで「逆翻訳」を練習し、AIが訳質を評価する語学学習アプリ。中国のSNS（小紅書）にて1,000件以上の反響を獲得。TypeScript / React / Docker / Edge TTS / VOICEVOX\\
\url{https://github.com/RainYaya/baku-yomi}

\textbf{AWS Mondai — AWS認定試験 学習プラットフォーム}\par
AWS認定13区分・4,591問を収録した日中英トリリンガル対応の学習サイト。自身の資格取得のために構築し、現在も稼働・利用者あり。TypeScript / Next.js / Vercel（ソースは非公開。デモは面接時に提示可能）

\textbf{設計上の判断（一例）}
\begin{itemize}
  \item 並行制御：GPU が1枚のため文字起こしは直列、ネットワーク律速のダウンロードは並列とし、競合するリソースによって方針を分けました。
  \item 状態管理：「再生時刻を唯一の状態源とする」方針を採用し、3つの表示面を同一の時計から導出。状態を二重に持たないことで同期のズレを構造的に排除しました。
  \item 鍵管理：BYOK 方式により API キーをサーバーへ送信しない設計とし、個人開発における責任範囲を限定しました。
\end{itemize}

% ============ 保有資格 ============
\kyariku{保有資格}

\begin{tabular}{@{}p{6em}p{\dimexpr\textwidth-6em\relax}@{}}
  2021年12月 & 日本語能力試験（JLPT）N1 合格 \\
  2026年07月 & AWS Certified Cloud Practitioner (CLF-C02) \\
  2026年07月 & AWS Certified Solutions Architect -- Associate (SAA-C03) \\
  2026年08月 & AWS Certified Solutions Architect -- Professional (SAP-C02) \\
\end{tabular}

% ============ 語学力 ============
\kyariku{語学力}

\begin{tabular}{@{}p{6em}p{\dimexpr\textwidth-6em\relax}@{}}
  中国語 & ネイティブ \\
  日本語 & JLPT N1（業務における会議・設計書作成・報告書作成が可能） \\
  英語 & 技術文書の読解が可能 \\
\end{tabular}

% ============ 自己PR ============
\kyariku{自己PR}

\textbf{【必要だと判断したことを、最短距離で習得する】}\par
現在参画しているクラウド移行案件では、インフラ基盤チームが AWS 側を担当する体制でしたが、開発側としても移行先の全体像を把握していなければ設計上の判断ができない場面がありました。そこで体系的な知識の獲得が必要だと判断し、AWS認定資格を CLF から SAP まで段階的に取得いたしました。実務経験としてはまだ浅いことを自覚しておりますが、基盤チームと同じ論点で議論できる土台は作れたと考えております。

\textbf{【曖昧な情報を、判断できる形に整理する】}\par
RPA 製品の UI 開発では、「操作が分かりにくい」という抽象的な要望に対し、実際の操作手順を確認してどの工程で迷いが生じているのかを特定し、具体的な改修に落とし込みました。現在の移行案件でも、ドキュメントに残っていない挙動について「仕様なのか不具合なのか」を有識者に確認しながら進めております。

また個人開発では、設計判断の理由と却下した代替案を ADR として文書化し、「後から読む人が理解できる状態」を意識した開発を心がけております。実装に生成AIを併用するようになってから、「何を作るか」「なぜそうするか」を自分で決めて言語化する必要がむしろ増えたと感じており、この習慣はチーム開発においても活きるものと考えております。

\vspace{8pt}
\begin{flushright}
以上
\end{flushright}

\end{document}
